\documentclass[a4paper, 12pt, french]{article}

% ==============
% IMPORT PACKAGE
% ==============

\usepackage[T1]{fontenc}
\usepackage[utf8]{inputenc}
\usepackage[margin=2cm]{geometry}
\usepackage{amsfonts, amssymb, amsmath}
\usepackage{lmodern}
\usepackage{theorem}
\usepackage{babel}

\date{\today}
\author{}
\title{RESOLUTION DE SYSTEME LINEAIRE}

\theoremstyle{break}
\theoremheaderfont{\bfseries}
\theorembodyfont{\upshape}

\newtheorem{defin}{Définition}[section]
\newtheorem{rema}[defin]{Remarque}
\newtheorem{theo}[defin]{Théoreme}
\newtheorem{algo}{Algorithme}

\renewcommand{\familydefault}{\sfdefault}

\begin{document}
\maketitle
Un système d'équation linéaire est un ensemble d'équation sous la forme :
$$\left\lbrace\begin{array}{lll}
    a_{11}x_{1} \: + \: a_{12}x_{2} \: + \: \cdots \: + \: a_{1n}x_{n} & = & b_{1} \\
    a_{21}x_{1} \: + \: a_{22}x_{2} \: + \: \cdots \: + \: a_{2n}x_{n} & = & b_{2} \\
    \vdots &  & \vdots \\
    a_{n1}x_{1} \: + \: a_{n2}x_{2} \: + \: \cdots \: + \: a_{nn}x_{n} & = & b_{n} \\
\end{array} \right .$$

\section{Forme matricielle}
Un système d'équation linéaire peut aussi s'écrire sous forme $Ax=b \; ou \; A \in M_{mn}(\mathbb{K}),\; x \in \mathbb{K}^{n} \; et \; b \in \mathbb{K}^{m}$
$$A = \left(\begin{array}{cccc}
    a_{11} & a_{12} & \cdots & a_{1n} \\
    a_{21} & a_{22} & \cdots & a_{2n} \\
    \vdots & \vdots & \ddots & \vdots \\
    a_{m1} & a_{m2} & \cdots & a_{mn} \\
    \end{array}\right) ; 
    x = \left(\begin{array}{cccc}
    x_{1}\\
    x_{2}\\
    \vdots \\
    x_{n} \\
    \end{array}\right)
    \; et \; 
    b = \left(\begin{array}{ccc}
    b_{1}\\
    b_{2}\\
    \vdots \\
    b_{n} \\
    \end{array}\right)$$


\begin{defin}
    Le rang d'une matrice notée $rg(A)$ est le nombre maximum de vecteurs ligne (ou colonne) linéairement indépendant.
\end{defin}
\begin{rema}
    Si $A \in M_{mn}(\mathbf{K}) \; alors \; rg(A) \leq \min(m,n)$
\end{rema}

\section{Résolution de système linéaire avec matrice carrée}
\subsection{Décomposition en LU}
On suppose que $A \in Gl_{n}(\mathbb{K})$ et on considère que le systeme linéaire $Ax=b$. L'idée est de décomposer $A=LU$ avec $\mathit{L}$ un matrice triangulaire inférieur et $\mathit{U}$ un matrice triangulaire supérieur.\\
Pour résoudre $Ax=b$ en 2 étape :
\begin{itemize}
    \item[-] Résoudre l'équation $Ly=b$ pour $y=Ux$
    \item[-] Puis résoudre $Ux=y$
\end{itemize}    
Il faut environ $\frac{2}{3}n^{3}$ pour calculer \textit{LU} et pour résoudre un système triangulaire il faut $n^2$ opération.\\
\textbf{Application}:\\
\textbf{Calcul de déterminant et inverse}:\\
Une fois obtenu la décomposition $LU$ on a $$A=LU \equiv detA\:=\:detL\:detU$$
Or si $detL=1\:(resp\:detU=1)$ on a : $$detA\:=\:detU=\prod_{i=1}^{n} u_{ii}(resp\:detA\:=\:detU=\prod_{i=1}^{n} l_{ii}).$$
Pour l'inverse on note $A^{-1}=(x_{1}|x_{2}| \cdots |x_{n})$ ou chacun des vecteurs $x_{i}$ est solution de $Ax_{i}=e_{i}$.\\ \\

\begin{algo}
	Entrée : A \\
	Sortie : L,U \\
	\begin{itemize}
	\item[-] \textbf{Si on suppose que les diagonales de L vaut 1}
	Composant de la matrice L :
	$$ l_{ij} = \frac{1}{u_{jj}}(a_{ij}-\sum_{k=1}^{j-1}l_{ik}u_{kj}) $$
	Composant de la matrice U :
	$$ u_{ij} = a_{ij} - \sum_{k=1}^{i-1}l_{ik}u_{kj}$$
	\item[-] \textbf{Si on suppose que les diagonales de U vaut 1}
	Composant de la matrice L :
	$$ l_{ij} = a_{ij} - \sum_{k=1}^{j-1}l_{ik}u_{kj} $$
	Composant de la matrice U :
	$$ u_{ij} = \frac{1}{l_{ii}}(a_{ij} - \sum_{k=1}^{i-1}l_{ik}u_{kj})$$
	\end{itemize}
\end{algo}

\section{Décomposition QR}
On suppose que $A \in Gl_{n}(\mathbb{K})$. 
L'idée est de décomposer $A=QR$ avec $\mathit{Q}$ un matrice orthogonale ($QQ^{t}=Q^{t}Q=I$) et $\mathit{R}$ un matrice triangulaire supérieur.
Pour pouvoir faciliter la résolution de $Ax=b$. On écrit $Ax=b \equiv QRx=b$, on multiplie membre à membre par $Q^{t}$ et on obtient $Rx=Q^{t}b$.\\
\begin{rema}
    Cette décomposition QR existe toujours.
\end{rema}
\begin{algo}
Entrée : A \\
Sortie : Q,R \\
Notons $[q_{1},q_{2},\cdots,q_{n}] \; et \; [a_{1},a_{2},\cdots,a_{n}]$ les vécteurs colonnes réspectifs des matrices Q et A.\\
Pour trouver Q : $$ Posons \; u_{1}=a_{1} \; ; \; q_{1}=\frac{u_{1}}{\|u_{1}\|}$$
Pour j=2,3,...,n: $$u_{j} = a_{j} - \sum_{i=1}^{j-1}<a_{j},q_{i}>q_{i} \; ; \; q_{j}=\frac{u_{j}}{\|u_{j}\|}$$
Pour trouver R, on résout l'équation $R=Q^{t}A$.
\end{algo}

\section{Méthode de Cholesky}
\begin{theo}
    Soit $A \in M_{n}(\mathbf{R})$ une matrice symétrique définie positive alors, il existe un unique $B \in M_{n}(\mathbf{R})$ triangulaire inférieur de coefficient diagonaux strictément positive tel que $A=BB^{t}$.\\
    Cette décomposition est appelée décompositionn de Cholesky.
\end{theo}
\begin{algo}
	Entrée : A \\
	Sortie : B \\
	Diagonal :  $$b_{jj}=\sqrt{a_{jj}-\sum_{k=1}^{j}b_{jk}^{2}}$$
	Non diagonaux : $$ b_{ij} = \frac{1}{b_{jj}}(a_{ij} - \sum_{k=1}^{j-1}b_{i}{k}b_{jk})$$
\end{algo}
\end{document}
